\section{Method}

\subsection{Forward Noising and Dense Prediction}

For a clean image $x_0$, timestep $t$, and Gaussian noise $\epsilon$, the forward
process is
\begin{equation}
  x_t = \alpha_t x_0 + \sigma_t \epsilon.
  \label{eq:forward}
\end{equation}
A conventional diffusion model predicts $\epsilon_\theta(x_t,t,c)$.
\method{} instead predicts the ordered component sum
in Eq.~\ref{eq:full-factorization}. The current implementation uses
unconditional models for the main MVP evidence. Class labels may exist in
datasets, but the restricted synthesis operator never receives labels, prompts,
timesteps, previous tokens, or image features. Figure~\ref{fig:method} shows the
dataflow and marks the blocked paths into $S_k$ that define the default method
boundary.

\begin{figure}[t]
  \centering
  \includegraphics[width=\linewidth]{method_overview.png}
  \caption{%
  \method{} factorizes dense pixel-space diffusion noise prediction into ordered
  denoising tokens. The predictor may be expressive, but each
  synthesis operator $S_k$ receives only the current token and is restricted to
  shallow, condition-free, bias-free local synthesis. Any prefix of components
  forms a valid partial noise prediction and therefore a valid pixel-space
  denoising result. Dashed blocked paths mark information that must not enter
  $S_k$, including $x_t$, timesteps, labels, prompts, previous tokens, skip
  features, and learned constants.}
  \label{fig:method}
\end{figure}

\subsection{Next Denoising Tokens}

The token predictor $T_k$ may be expressive. Our default implementation sees
the noised image, timestep, and previous tokens:
\begin{align}
  z^{(1)} &= T_1(x_t,t), \\
  z^{(k)} &= T_k(x_t,z^{(<k)},t), \quad k>1.
\end{align}
In the tiny MVP model, $T_k$ is a small U-Net-like predictor with token heads.
Autoregressive feedback is not part of the synthesis restriction or the core
factorization claim: a simultaneous variant may instead use
$z^{(k)}=T_k(x_t,t)$. We evaluate this choice directly.

\subsection{Restricted Synthesis Operator}

Each token is expanded independently:
\begin{equation}
  \epsilon^{(k)} = S_k\left(z^{(k)}\right).
\end{equation}
The default $S_k$ is a bias-free projection followed by a bias-free local
convolution. It is token-only and condition-free. Because no layer has bias and
there is no learned constant input, zero tokens map to zero components:
\begin{equation}
  S_k(0)=0.
\end{equation}
This restriction separates \method{} from a learned decoder. A deep-$S_k$
ablation is included only to test degeneration risk: it uses biased nonlinear
convolutions while preserving the token-only interface. This ablation can
improve apparent endpoint metrics but fails the zero-token diagnostic, so it is
not the default method.

\subsection{Prefix Denoising}

For any prefix length $m$, \method{} forms the prefix prediction and pixel-space
denoising result in Eq.~\ref{eq:prefix-denoise}. Thus every prefix is a valid
partial noise prediction and every prefix produces a pixel-space denoising
result. This is the primary interface exposed by the method.

\subsection{Denoise-Path Objective}

The current MVP objective is
\begin{equation}
  \mathcal{L}
  =
  \mathcal{L}_{\epsilon}
  + 0.15\,\mathcal{L}_{\mathrm{path\ prefix}}
  + 0.30\,\mathcal{L}_{\mathrm{path\ component}}.
\end{equation}
The denoise path interpolates from the zero-update $x_t$ reconstruction toward
coarse-to-fine spatial targets. The progress schedule uses power 1.5 in the
current best MVP configuration. This objective deliberately avoids forcing early
prefixes to be immediately clean; instead, prefixes are trained to follow a
partial denoising path.
